\chapter{Information Theory: Entropy, Cross-Entropy, and KL Divergence}
\label{chap:ch07}

% Chapter goal: measure uncertainty and the mismatch between
% distributions, and show that the cross-entropy loss and the KL
% divergence behind VAEs, distillation, and RLHF are the same idea.

\Cref{chap:ch06} derived the cross-entropy loss from maximum
likelihood: a classifier that is trained to make the observed labels
probable ends up minimizing $-\sum_k y_k \log \hat{p}_k$. The derivation
was complete, but it left the name unexplained. ``Entropy'' is a word
from physics, and it entered machine learning through \emph{information
theory}, the mathematics of communication that Claude
Shannon~\cite{shannon1948mathematical} created in 1948 to answer a
practical question: how many bits does it take to send a message?

The answer turns out to depend on how predictable the messages are, and
the quantity that measures this, entropy, is a property of a probability
distribution. This chapter builds three related quantities from it. The
\emph{entropy} of a distribution measures how uncertain it is. The
\emph{cross-entropy} measures the cost of describing data from one
distribution with a model of another. And the \emph{KL divergence}
measures how far apart two distributions are. Together they give a
second, independent reading of almost every loss function in the book.

\section{Entropy}
\label{sec:ch07-entropy}

\subsection{Surprise}
\label{sec:ch07-surprise}

Some news tells you more than other news. If a weather forecast for the
desert says ``no rain tomorrow'', you have learned almost nothing; you
expected that. If it says ``snow tomorrow'', you have learned a great
deal. The less probable an event, the more \emph{surprising} it is when
it happens, and the more information its occurrence carries. Information
theory starts by turning this intuition into a number.

We want a function $s(p)$ that assigns a surprise to an event of
probability $p$. Three requirements pin it down.
\begin{itemize}
  \item \emph{Rarer is more surprising.} $s(p)$ should decrease as $p$
    grows.
  \item \emph{Certain events are no surprise.} An event of probability
    $1$ tells us nothing, so $s(1) = 0$.
  \item \emph{Independent surprises add.} If two independent events of
    probabilities $p$ and $q$ both happen, the probability of the pair is
    $pq$ (\Cref{sec:ch05-independence}). Learning both should be as
    informative as learning one and then the other, so $s(pq) = s(p) +
    s(q)$.
\end{itemize}
The last requirement asks for a function that turns products into sums,
and that is what the logarithm does. Up to the choice of base, the only
function meeting all three is $s(p) = -\log p$. The minus sign makes it
positive, since $\log p \le 0$ for $p \le 1$.

\begin{definition}[Surprise]
  \label{def:surprise}
  The \emph{surprise}, or \emph{information content}, of an outcome $x$
  of probability $p(x) > 0$ is
  \[
    s(x) = -\log p(x) = \log \frac{1}{p(x)} .
  \]
  With $\log_2$ it is measured in \emph{bits}; with the natural logarithm
  it is measured in \emph{nats}.
\end{definition}

The base only sets the unit, as meters and feet do for length:
$\log_2 p = \log p / \log 2$, so one nat is $1/\log 2 \approx 1.44$
bits. Bits are easier to picture. A fair coin landing heads has
probability $\tfrac12$ and surprise $-\log_2 \tfrac12 = 1$ bit, the
information in one yes-or-no answer. An outcome of probability
$\tfrac18$ has surprise $3$ bits: it singles out one of eight equally
likely cases, as three yes-or-no answers do. Each halving of the
probability adds one bit, and an event of probability $0.01$ carries
$\log_2 100 \approx 6.64$ bits. Machine learning code uses the natural
logarithm, and so, unless we say otherwise, does this book; we switch to
bits when the counting interpretation helps.

The negative log-probability is not new. It is exactly the term that
each data point contributes to the negative log-likelihood of
\Cref{sec:ch06-mle-flow}. Minimizing the negative log-likelihood
therefore means making the training data as unsurprising to the model as
possible, a reading that \Cref{sec:ch07-entropy-ai} and the rest of the
chapter develop.

\subsection{Entropy: Average Surprise}
\label{sec:ch07-entropy-def}

A single outcome has a surprise. A random variable has many possible
outcomes, each with its own surprise, and the natural summary is the
average, weighted by how often each outcome occurs.

\begin{definition}[Entropy]
  \label{def:entropy}
  The \emph{entropy} of a discrete random variable $X$ with PMF $p$ is
  its expected surprise,
  \[
    H(X) = \E\bigl[-\log p(X)\bigr] = -\sum_{x} p(x) \log p(x) .
  \]
  Terms with $p(x) = 0$ are taken to be $0$. We also write $H(p)$, or
  $H(\vect{p})$ for a probability vector $\vect{p}$, when it is the
  distribution rather than the random variable that matters.
\end{definition}

The convention for zero probabilities is the natural one. An outcome
that never occurs contributes nothing to the average, and indeed $p \log
p \to 0$ as $p \to 0$: the surprise $-\log p$ grows without bound, but
its weight $p$ shrinks faster. Entropy depends only on the probabilities,
not on the values $x$ themselves; relabeling the outcomes does not change
it.

\begin{example}[Coins and dice]
  \label{ex:ch07-coins}
  A fair coin has entropy
  \[
    H = -\tfrac12 \log_2 \tfrac12 - \tfrac12 \log_2 \tfrac12 = 1 \text{ bit}.
  \]
  A coin with probability of heads $0.9$ is more predictable. Heads carries
  $-\log_2 0.9 \approx 0.15$ bits and tails $-\log_2 0.1 \approx 3.32$
  bits, but tails is rare, so the average is
  \[
    H = 0.9 \times 0.15 + 0.1 \times 3.32 \approx 0.47 \text{ bits}.
  \]
  A coin that always lands heads has $H = -1 \cdot \log_2 1 = 0$: there
  is nothing to learn from flipping it. A fair six-sided die has six
  outcomes of surprise $\log_2 6$ each, so $H = \log_2 6 \approx 2.58$
  bits, more than a coin, because there are more ways for it to land.
\end{example}

For a Bernoulli variable with parameter $\theta$, the entropy is the
\emph{binary entropy function}
\[
  H(\theta) = -\theta \log \theta - (1 - \theta) \log (1 - \theta),
\]
plotted on the right of \Cref{fig:ch07-entropy}. It is $0$ at both ends,
where the outcome is certain, rises steeply away from them, and peaks at
$\theta = \tfrac12$, where the coin is hardest to predict. The curve is
flat near the top: a coin with $\theta = 0.7$ still has $0.88$ bits of
entropy. It is steep near the ends, so moving from $\theta = 0.99$ to
$0.9$ multiplies the entropy by about six, from $0.08$ to $0.47$ bits.

\begin{figure}[ht]
  \centering
  \begin{tikzpicture}
    \begin{axis}[
      name=surp, width=0.48\textwidth, height=0.36\textwidth,
      xmin=0, xmax=1, ymin=0, ymax=7,
      xtick={0, 0.125, 0.25, 0.5, 1},
      xticklabels={$0$, , $\tfrac14$, $\tfrac12$, $1$},
      ytick={0, 1, 2, 3, 6},
      xlabel={$p$}, ylabel={$-\log_2 p$ (bits)},
      axis lines*=left, title={\small surprise},
    ]
      \addplot[very thick, blue!60!black, domain=0.008:1, samples=200]
        {-ln(x)/ln(2)};
      \addplot[only marks, mark=*, mark size=1.6pt, blue!60!black]
        coordinates {(0.5,1) (0.25,2) (0.125,3)};
    \end{axis}
    \begin{axis}[
      at={(surp.east)}, anchor=west, xshift=0.07\textwidth,
      width=0.48\textwidth, height=0.36\textwidth,
      xmin=0, xmax=1, ymin=0, ymax=1.1,
      xtick={0, 0.5, 0.9, 1}, ytick={0, 0.47, 1},
      xlabel={$\theta$}, ylabel={$H(\theta)$ (bits)},
      axis lines*=left, title={\small binary entropy},
    ]
      \addplot[very thick, blue!60!black, domain=0.0005:0.9995, samples=200]
        {-(x*ln(x) + (1-x)*ln(1-x))/ln(2)};
      \addplot[dashed, accent] coordinates {(0.9,0) (0.9,0.469)};
      \fill[accent] (axis cs:0.9,0.469) circle (2pt);
      \fill[accent] (axis cs:0.5,1) circle (2pt);
    \end{axis}
  \end{tikzpicture}
  \caption{Left: the surprise of an outcome of probability $p$. Each
    halving of $p$ adds one bit (dots at $p = \tfrac12, \tfrac14,
    \tfrac18$), and it grows without bound as $p \to 0$. Right: the
    entropy of a coin with probability of heads $\theta$. It is largest,
    $1$ bit, for a fair coin, and $0.47$ bits for $\theta = 0.9$ (red).}
  \label{fig:ch07-entropy}
\end{figure}

\subsection{Entropy as the Length of a Description}
\label{sec:ch07-coding}

Shannon's interpretation of entropy is as a count of yes-or-no questions,
or equivalently of bits in a message. A weather station on a remote
mountain reports once a day, over an expensive satellite link, one of
four conditions:
\begin{center}
  \begin{tabular}{@{}lcccc@{}}
    \toprule
    weather & sun & cloud & rain & snow \\
    \midrule
    probability & $\tfrac12$ & $\tfrac14$ & $\tfrac18$ & $\tfrac18$ \\
    code & \texttt{0} & \texttt{10} & \texttt{110} & \texttt{111} \\
    \bottomrule
  \end{tabular}
\end{center}
The obvious encoding uses two bits per report, \texttt{00}, \texttt{01},
\texttt{10}, \texttt{11}, one for each condition. The code in the last row
does better by giving short codewords to common conditions and long ones
to rare conditions. No codeword is the beginning of another, so a string
of reports such as \texttt{0110010} can be split up without separators:
\texttt{0}, \texttt{110}, \texttt{0}, \texttt{10}, that is, sun, rain,
sun, cloud. The average length is
\[
  \tfrac12 \cdot 1 + \tfrac14 \cdot 2 + \tfrac18 \cdot 3 + \tfrac18 \cdot 3
  = 1.75 \text{ bits},
\]
a saving of $12.5\%$. Each codeword is the sequence of answers to a game
of twenty questions: ``Is it sunny?'' If not, ``Is it cloudy?'' If not,
``Is it raining?''

Now compute the entropy of the weather. The surprises are $1$, $2$, $3$,
and $3$ bits, exactly the codeword lengths, so
\[
  H = \tfrac12 \cdot 1 + \tfrac14 \cdot 2 + \tfrac18 \cdot 3 + \tfrac18
  \cdot 3 = 1.75 \text{ bits}.
\]
This is no coincidence. Shannon's \emph{source coding theorem} says that
no code of this kind can use fewer than $H(X)$ bits per report on
average, and that codes using close to $H(X)$ bits exist for every
distribution. The ideal is to spend $-\log_2 p(x)$ bits on outcome $x$:
rare outcomes cost more, common ones less. When those lengths are whole
numbers, as here, the ideal is reached exactly; otherwise, encoding long
blocks of reports at once gets as close to $H$ bits per report as we
like. So entropy is the length of the shortest possible description of
$X$: the number of bits that \emph{must} be sent to say which outcome
occurred. A predictable source needs few bits, an unpredictable one many,
and a uniform choice among $K$ outcomes needs $\log_2 K$.

This view also explains the ``cross'' in cross-entropy, the subject of
the next section. A code built for the wrong probabilities, for example
one designed in summer and used in winter, spends more bits than
necessary, and the cross-entropy is how many.

\subsection{Properties of Entropy}
\label{sec:ch07-entropy-props}

The examples suggest that entropy lies between two extremes: $0$ when the
outcome is certain and $\log K$ when all $K$ outcomes are equally likely.
The proof rests on one inequality,
\begin{equation}
  \log t \le t - 1 \quad \text{for all } t > 0,
  \text{ with equality only at } t = 1 .
  \label{eq:ch07-log-ineq}
\end{equation}
The right side is the tangent line of $\log t$ at $t = 1$
(\Cref{sec:ch02-linear-approx}), and a concave function lies below its
tangent lines (\Cref{fig:ch01-concave}).

\begin{proposition}[Range of entropy]
  \label{prop:ch07-entropy-range}
  Let $X$ take at most $K$ values. Then
  \[
    0 \le H(X) \le \log K .
  \]
  $H(X) = 0$ exactly when one outcome has probability $1$, and $H(X) =
  \log K$ exactly when all $K$ outcomes have probability $1/K$.
\end{proposition}

\begin{proof}
  Each term $-p(x) \log p(x)$ is at least $0$, because $0 \le p(x) \le 1$
  makes $\log p(x) \le 0$. A term is $0$ only if $p(x)$ is $0$ or $1$, so
  the sum is $0$ only if every probability is $0$ or $1$, that is, if one
  outcome is certain.

  For the upper bound, sum over the outcomes with $p(x) > 0$ and apply
  \cref{eq:ch07-log-ineq} with $t = 1/(K p(x))$:
  \[
    H(X) - \log K = \sum_x p(x) \log \frac{1}{K p(x)}
    \le \sum_x p(x) \Bigl( \frac{1}{K p(x)} - 1 \Bigr)
    = \sum_x \frac{1}{K} - 1 \le 0 .
  \]
  The first step uses $\sum_x p(x) = 1$ to write $\log K$ as $\sum_x p(x)
  \log K$. The last sum has at most $K$ terms, so it is at most $1$.
  Equality requires $t = 1$, that is, $p(x) = 1/K$, for every outcome.
\end{proof}

Entropy is largest for the uniform distribution, so the uniform is the
``most uncertain'' or ``least informative'' distribution on $K$ outcomes.
Any knowledge that makes some outcomes more likely than others lowers the
entropy. This is the basis of the \emph{principle of maximum entropy}:
when we must choose a distribution and know only a few facts about it,
choose the one with the largest entropy among those consistent with the
facts, since every other choice assumes something we do not know.

Two more properties are used constantly.

\paragraph{Entropy is concave.}
Viewed as a function on the probability simplex $\Delta_K$ of
\Cref{sec:ch01-domain-range}, $H(\vect{p}) = \sum_k \phi(p_k)$ with
$\phi(p) = -p \log p$. Since $\phi''(p) = -1/p < 0$, each term is concave,
and a sum of concave functions is concave (\Cref{def:concave}). In words,
mixing two distributions never produces less entropy than the average of
their entropies. Pick one of two coins at random, a fair one with
entropy $1$ bit and a two-headed one with entropy $0$, and flip it. The
flip lands heads with probability $\tfrac34$, so its entropy is
$H(\tfrac34) \approx 0.81$ bits, more than the average $0.5$: not knowing
which coin was picked adds uncertainty of its own.
Concavity also makes maximum entropy problems easy, because maximizing a
concave function is minimizing a convex one, and every local optimum is
global (\Cref{thm:ch01-local-global}).

\paragraph{Entropies of independent variables add.}
The \emph{joint entropy} of two random variables is the entropy of the
pair, $H(X, Y) = -\sum_{x, y} p(x, y) \log p(x, y)$. If $X$ and $Y$ are
independent, then $p(x, y) = p(x)\,p(y)$, the log splits into $\log p(x) +
\log p(y)$, and
\[
  H(X, Y) = H(X) + H(Y) .
\]
This is the additivity of surprise from \Cref{sec:ch07-surprise}, now
averaged. Flipping a fair coin and rolling a fair die gives one of $12$
equally likely pairs, with entropy $\log_2 12 \approx 3.58$ bits, which is
$1$ bit for the coin plus $2.58$ for the die. For $N$ i.i.d.\ draws, the
entropy is $N$ times that of one draw. When $X$ and $Y$ are dependent,
knowing one tells us something about the other, and $H(X, Y)$ is smaller
than the sum; the amount of the shortfall, called the \emph{mutual
information}, measures how much they share.

\subsection{Continuous Variables: Differential Entropy}
\label{sec:ch07-differential}

For a continuous random variable with density $f$, replacing the sum by
an integral gives the \emph{differential entropy}
\[
  h(X) = -\int_{-\infty}^{\infty} f(x) \log f(x)\, dx .
\]
It behaves like entropy in most ways, but it is not a count of bits. A
uniform variable on $[a, b]$ has density $1/(b - a)$ there, so
\[
  h(X) = -\int_a^b \frac{1}{b - a} \log \frac{1}{b - a}\, dx = \log(b - a) .
\]
Wider intervals have larger entropy, as expected, but the interval
$[0, \tfrac12]$ has $h = \log \tfrac12 < 0$. A negative entropy cannot be
a number of questions. The reason is that pinning down a real number
exactly takes infinitely many bits; differential entropy drops this
infinite part and keeps only what depends on the distribution. It also
depends on the units: measuring the same quantity in millimeters
instead of meters stretches the density by $1000$ and adds $\log 1000$ to
$h$. Differences of differential entropies, and the KL divergence later
in this chapter, do not have these problems.

\begin{example}[Gaussian]
  \label{ex:ch07-gaussian-entropy}
  For $X \sim \mathcal{N}(\mu, \sigma^2)$, the log-density is
  $-\tfrac12 \log(2\pi\sigma^2) - (x - \mu)^2/(2\sigma^2)$, so
  \[
    h(X) = \E\bigl[-\log f(X)\bigr]
    = \tfrac12 \log(2\pi\sigma^2) + \frac{\E[(X - \mu)^2]}{2\sigma^2}
    = \tfrac12 \log(2\pi\sigma^2) + \tfrac12
    = \tfrac12 \log(2\pi e \sigma^2) .
  \]
  The mean does not appear: shifting a distribution does not change its
  uncertainty. The standard Gaussian has $h \approx 1.42$ nats, and
  doubling $\sigma$ adds $\log 2$, exactly one bit. For a multivariate
  Gaussian $\mathcal{N}(\vect{\mu}, \Sigma)$ in $\R^d$ the same calculation
  gives
  \[
    h = \tfrac12 \log \det(2\pi e\, \Sigma) ,
  \]
  in which the determinant of \Cref{sec:ch03-det} measures the volume
  of the cloud of likely points.
\end{example}

The Gaussian has a special place here. Among all densities on $\R$ with a
given variance $\sigma^2$, the Gaussian $\mathcal{N}(\mu, \sigma^2)$ has
the largest differential entropy. By the principle of maximum entropy,
if all we know about a noise source is its variance, the Gaussian is the
assumption that adds nothing else. This is one more reason, beside the
central limit theorem, that Gaussian noise is the default in the
regression models of \Cref{sec:ch06-mle-ml}. On the interval $[a, b]$
with no other constraint, the maximum entropy density is the uniform
one, mirroring \Cref{prop:ch07-entropy-range}.

\subsection{Entropy in AI}
\label{sec:ch07-entropy-ai}

A classifier's output is a categorical distribution, and its entropy
measures how unsure the classifier is. With four classes, compare
\[
  \hat{\vect{p}} = (0.97, 0.01, 0.01, 0.01),
  \qquad H \approx 0.24 \text{ bits},
  \qquad\qquad
  \hat{\vect{p}} = (0.4, 0.3, 0.2, 0.1),
  \qquad H \approx 1.85 \text{ bits}.
\]
The first prediction is confident and the second is close to the maximum
of $\log_2 4 = 2$ bits, a shrug. Unlike the top probability alone,
entropy accounts for how the rest of the mass is spread. Predictions with
high entropy are the ones worth a second look: they can be routed to a
human, or, in \emph{active learning}, sent to be labeled first, because
the model expects to learn most from them.

The temperature of \Cref{fig:ch01-softmax} is an entropy dial. For the
logits $\vect{z} = (2, 1, 0.5, -1)$, the softmax has entropy $0.78$ bits
at $T = 0.5$, $1.46$ bits at $T = 1$, and $1.97$ bits at $T = 5$, close
to the maximum of $2$. Lowering the temperature of a language model makes
each next-token distribution lower in entropy, and the text more
predictable; raising it does the opposite.

\begin{aibox}
  Entropy appears wherever a model needs to measure, reward, or limit
  uncertainty.
  \begin{itemize}
    \item \emph{Decision trees} choose each split to reduce the entropy of
      the labels. Ten emails, five of them spam, have label entropy $1$
      bit. If a split sends four spam and one non-spam email left, and one
      spam and four non-spam right, each side has entropy $H(0.8) \approx
      0.72$ bits, so the split has an \emph{information gain} of $1 - 0.72
      = 0.28$ bits.
    \item \emph{Reinforcement learning} often adds an entropy bonus to
      the objective of a policy, which keeps its action distribution from
      collapsing onto one action too early and so keeps it exploring.
      Maximum entropy methods such as soft
      actor-critic~\cite{haarnoja2018sac} build this into the objective.
    \item \emph{Perplexity}, the standard score of a language model, is
      $e$ raised to its average negative log-probability per token
      (\Cref{sec:ch06-mle-ml}). For a model whose predictions match the
      data, this is $e^{H}$, the \emph{effective number of choices}: a
      uniform distribution over $K$ tokens has perplexity exactly $K$, and
      a perplexity of $20$ means the model is, on average, as uncertain as
      a fair choice among $20$ tokens. The next section makes this
      precise for models that do not match the data.
  \end{itemize}
\end{aibox}

\subsection{In Code: Computing Entropy}
\label{sec:ch07-entropy-code}

The definition translates directly into NumPy~\cite{harris2020numpy}.
For the weather station of \Cref{sec:ch07-coding}:

\begin{lstlisting}
import numpy as np

p = np.array([0.5, 0.25, 0.125, 0.125])
-np.sum(p * np.log2(p))           # 1.75 bits
-np.sum(p * np.log(p))            # 1.213 nats
\end{lstlisting}

The direct formula breaks as soon as a probability is zero, because
\texttt{np.log(0)} is $-\infty$ and $0 \times (-\infty)$ is
\texttt{nan}. SciPy~\cite{virtanen2020scipy} provides versions that apply
the convention $0 \log 0 = 0$:

\begin{lstlisting}
from scipy import special, stats

q = np.array([0.7, 0.3, 0.0])
-np.sum(q * np.log(q))            # nan, with a warning
special.entr(q).sum()             # 0.611: entr(x) = -x log x, entr(0) = 0
stats.entropy(q)                  # 0.611 nats
stats.entropy(p, base=2)          # 1.75 bits
stats.entropy([5, 3, 2], base=2)  # 1.485: counts are normalized first
\end{lstlisting}

\begin{warning}
  \texttt{stats.entropy} quietly divides its input by the sum, so a vector
  that does not sum to $1$ gives a number rather than an error. Given a
  second argument, \texttt{stats.entropy(p, q)} computes the KL
  divergence, which comes later in this chapter, not an entropy. In
  PyTorch, \texttt{torch.distributions.Categorical(probs=p).entropy()}
  computes entropy in nats.
\end{warning}

The normalization makes it easy to estimate the entropy of a distribution
from samples: count how often each outcome occurs and take the entropy of
the counts. With plenty of data, the estimate is close to the truth. With
little data it is too low on average, because a small sample looks less
varied than the distribution it came from:

\begin{lstlisting}
rng = np.random.default_rng(0)
x = rng.choice(4, size=10_000, p=p)
counts = np.bincount(x, minlength=4)          # [4990 2505 1300 1205]
stats.entropy(counts, base=2)                 # 1.751, true value 1.75

estimates = [stats.entropy(np.bincount(rng.choice(4, size=20, p=p),
                                       minlength=4), base=2)
             for _ in range(1000)]
np.mean(estimates)                            # 1.630: biased low
\end{lstlisting}

The temperature dial of \Cref{sec:ch07-entropy-ai} takes one loop, using
SciPy's softmax:

\begin{lstlisting}
logits = np.array([2.0, 1.0, 0.5, -1.0])
for T in [0.1, 0.5, 1.0, 5.0, 100.0]:
    print(T, stats.entropy(special.softmax(logits / T), base=2))
# 0.1  0.0007   nearly one-hot
# 0.5  0.776
# 1.0  1.463
# 5.0  1.968
# 100  2.000    nearly uniform: log2(4) = 2
\end{lstlisting}

Finally, every continuous distribution in \texttt{scipy.stats} has an
\texttt{entropy} method that returns its differential entropy in nats,
which confirms \Cref{ex:ch07-gaussian-entropy} and the negative entropy
of a narrow uniform:

\begin{lstlisting}
stats.norm(scale=1.0).entropy()       # 1.419 = 0.5 * log(2 pi e)
stats.norm(scale=2.0).entropy()       # 2.112: log(2) = 0.693 more
stats.uniform(scale=0.5).entropy()    # -0.693 = log(1/2)
\end{lstlisting}

\section{Cross-Entropy}
\label{sec:ch07-cross-entropy}

\subsection{Coding with the Wrong Distribution}
\label{sec:ch07-ce-idea}

The weather station of \Cref{sec:ch07-coding} was designed in summer,
when the weather was sun, cloud, rain, or snow with probabilities
$\tfrac12$, $\tfrac14$, $\tfrac18$, $\tfrac18$, and its code gave sun the
shortest codeword. Winter comes, and the probabilities change to
$\tfrac18$, $\tfrac14$, $\tfrac18$, $\tfrac12$: snow is now the common
report. Nobody updates the code, so snow still costs three bits and sun
one:
\begin{center}
  \begin{tabular}{@{}lcccc@{}}
    \toprule
    weather & sun & cloud & rain & snow \\
    \midrule
    summer probability & $\tfrac12$ & $\tfrac14$ & $\tfrac18$ & $\tfrac18$ \\
    winter probability & $\tfrac18$ & $\tfrac14$ & $\tfrac18$ & $\tfrac12$ \\
    code (built for summer) & \texttt{0} & \texttt{10} & \texttt{110}
      & \texttt{111} \\
    \bottomrule
  \end{tabular}
\end{center}
The average length of a winter report is now
\[
  \tfrac18 \cdot 1 + \tfrac14 \cdot 2 + \tfrac18 \cdot 3 + \tfrac12 \cdot 3
  = 2.5 \text{ bits},
\]
worse than the plain two-bit code. The winter weather is no more
uncertain than the summer weather, since its probabilities are the same
numbers in a different order, and its entropy is still $1.75$ bits. A
code built for winter would need only that. The extra $0.75$ bits per
report are the price of using a code designed for the wrong
distribution.

In general, a code built for a distribution $q$ spends $-\log q(x)$ bits
on outcome $x$, the surprise that $q$ assigns to $x$. When the outcomes
actually come from $p$, the average cost is the surprise under $q$,
averaged under $p$. That is the cross-entropy.

\subsection{Definition and Properties}
\label{sec:ch07-ce-def}

\begin{definition}[Cross-entropy]
  \label{def:cross-entropy}
  The \emph{cross-entropy} of a distribution $q$ relative to a
  distribution $p$, on the same discrete outcomes, is
  \[
    H(p, q) = \E_{X \sim p}\bigl[-\log q(X)\bigr]
    = -\sum_x p(x) \log q(x) .
  \]
  Terms with $p(x) = 0$ are taken to be $0$. If $q(x) = 0$ for some $x$
  with $p(x) > 0$, then $H(p, q) = \infty$.
\end{definition}

The two arguments play different roles. The first, $p$, is the truth:
it decides how often each outcome occurs, and so how the surprises are
weighted. The second, $q$, is the model: it decides how surprising each
outcome is. The subscript in $\E_{X \sim p}$ makes this explicit; it says
that $X$ is drawn from $p$. When the two are the same, the cross-entropy
is the entropy, $H(p, p) = H(p)$. The weather station shows that using
the wrong $q$ costs extra, and this is true in general.

\begin{proposition}[Gibbs' inequality]
  \label{prop:ch07-gibbs}
  For any distributions $p$ and $q$ on the same outcomes,
  \[
    H(p, q) \ge H(p),
  \]
  with equality only when $q = p$.
\end{proposition}

\begin{proof}
  If $H(p, q) = \infty$ there is nothing to prove, so assume $q(x) > 0$
  wherever $p(x) > 0$. Sum over those outcomes and apply
  \cref{eq:ch07-log-ineq} with $t = q(x)/p(x)$:
  \[
    H(p) - H(p, q) = \sum_x p(x) \log \frac{q(x)}{p(x)}
    \le \sum_x p(x) \Bigl( \frac{q(x)}{p(x)} - 1 \Bigr)
    = \sum_x q(x) - 1 \le 0 .
  \]
  The last sum runs over only some outcomes, so it is at most $1$.
  Equality in the middle step requires $q(x)/p(x) = 1$ for every outcome
  with $p(x) > 0$, and then $q$ has no probability left for the others.
\end{proof}

\Cref{prop:ch07-entropy-range} is the special case where $q$ is uniform:
$H(p, q) = \log K$ for every $p$, so $\log K \ge H(p)$. In words, the
best description of data from $p$ is one built for $p$ itself, and the
entropy $H(p)$ is the floor that no model can go below. The excess
$H(p, q) - H(p)$, the $0.75$ bits of the weather station, is the
subject of the next section.

Cross-entropy is not symmetric: in general, $H(p, q) \ne H(q, p)$. With
the summer distribution $p_{\text{s}} = (\tfrac12, \tfrac14, \tfrac18,
\tfrac18)$ and the uniform distribution $u$ on four outcomes,
\[
  H(u, p_{\text{s}}) = \tfrac14 (1 + 2 + 3 + 3) = 2.25 \text{ bits},
  \qquad
  H(p_{\text{s}}, u) = 2 \text{ bits}.
\]
Using the summer code on uniformly random weather costs $2.25$ bits;
using the plain two-bit code on summer weather costs $2$. The asymmetry
is strongest when the model rules things out. If $q(x) = 0$ for an
outcome that does occur, the cross-entropy is infinite, because the
model was infinitely surprised; the reverse mistake, giving probability
to an outcome that never occurs, costs only a finite amount. A model
trained with cross-entropy therefore learns never to be certain that
something is impossible.

For continuous distributions with densities $p$ and $q$, the sum becomes
an integral, $H(p, q) = -\int p(x) \log q(x)\, dx$, and Gibbs' inequality
holds with the differential entropy $h(p)$ in place of $H(p)$.

\subsection{Cross-Entropy as a Loss Function}
\label{sec:ch07-ce-loss}

We can now explain the name of the loss in \Cref{sec:ch06-mle-ml}. A
classifier sees an input $\vect{x}$ and outputs a distribution $\vect{q} =
\operatorname{softmax}(\vect{z})$ over $K$ classes. The truth for this
input is the label, which, written as a one-hot vector $\vect{y}$
(\Cref{sec:ch05-categorical}), is a distribution that puts all its mass
on the correct class $c$. The cross-entropy of the prediction relative
to the label has only one nonzero term:
\[
  H(\vect{y}, \vect{q}) = -\sum_{k=1}^{K} y_k \log q_k = -\log q_c .
\]
This is the surprise of the model at the correct answer, and it is
exactly the negative log-likelihood of one labeled example under the
categorical model. Averaged over a data set, the cross-entropy loss and
the negative log-likelihood are the same number:
\[
  \frac{1}{N} \sum_{i=1}^{N} H(\vect{y}_i, \vect{q}_i)
  = -\frac{1}{N} \sum_{i=1}^{N} \log q_{i, c_i}
  = -\frac{1}{N}\, \ell(\vect{w}) .
\]
Training by minimizing the cross-entropy and training by maximum
likelihood are one and the same. Information theory reads the loss as
the average number of nats the model would need to encode the correct
labels; statistics reads it as how improbable the model finds them.

For the logits $\vect{z} = (2, 1, 0.5, -1)$ of \Cref{fig:ch01-softmax},
the prediction is $\vect{q} \approx (0.61, 0.22, 0.14, 0.03)$. If the
correct class is the first, the loss is $-\log 0.61 \approx 0.50$ nats;
if it is the second, $-\log 0.22 \approx 1.50$; if it is the last,
$-\log 0.03 \approx 3.50$. The loss is $0$ only for a prediction of
exactly $1$ on the correct class, and it grows without bound as that
probability goes to $0$, as the left of \Cref{fig:ch07-ce-vs-se} shows.

With two classes, $\vect{y} = (y, 1 - y)$ and $\vect{q} = (\hat{p}, 1 -
\hat{p})$, and the cross-entropy is
\[
  H(\vect{y}, \vect{q}) = -y \log \hat{p} - (1 - y) \log (1 - \hat{p}),
\]
the binary cross-entropy of logistic regression. The binary entropy of
\Cref{fig:ch07-entropy} is the special case $\hat{p} = y$.

\paragraph{The gradient.}
The loss is usually computed from the logits, and its gradient with
respect to them is remarkably simple. Writing out the softmax,
\[
  -\log q_c = -\log \frac{e^{z_c}}{\sum_j e^{z_j}}
  = -z_c + \log \sum_{j=1}^{K} e^{z_j} .
\]
The derivative of $-z_c$ with respect to $z_k$ is $-1$ if $k = c$ and $0$
otherwise, that is, $-y_k$. The derivative of the second term is, by the
chain rule, $e^{z_k} / \sum_j e^{z_j} = q_k$. So
\begin{equation}
  \frac{\partial}{\partial z_k} H(\vect{y}, \vect{q}) = q_k - y_k,
  \qquad
  \nabla_{\vect{z}} H(\vect{y}, \vect{q}) = \vect{q} - \vect{y} .
  \label{eq:ch07-ce-grad}
\end{equation}
Prediction minus target: the same form as the gradient $\hat{p} - y$ of
logistic regression in \Cref{sec:ch06-mle-ml}. For the logits above, with
the first class correct, the gradient is $(0.61 - 1, 0.22, 0.14, 0.03) =
(-0.39, 0.22, 0.14, 0.03)$. Gradient descent raises the correct logit and
lowers the others, each in proportion to the probability it wrongly
holds. The second term, $\log \sum_j e^{z_j}$, is called the
\emph{log-sum-exp} of the logits; it is convex, so the loss is convex in
the logits.

\paragraph{Why not squared error?}
Probabilities are numbers, so one could train a classifier by minimizing
the squared distance $\|\vect{q} - \vect{y}\|^2$ instead. Cross-entropy is
preferred for two reasons. The first is the one of \Cref{chap:ch06}: it
is the negative log-likelihood of the correct model of a label, while
squared error assumes Gaussian noise, which labels do not have. The
second is about gradients. Take a binary classifier with $\hat{p} =
\sigma(z)$ and a positive example, $y = 1$, which the model gets badly
wrong, with $z = -5$ and $\hat{p} \approx 0.007$. The cross-entropy has
gradient $\hat{p} - 1 \approx -0.99$ with respect to $z$, a strong push
in the right direction. The squared error $(1 - \hat{p})^2$ has gradient
$-2(1 - \hat{p})\,\sigma'(z) = -2(1 - \hat{p})^2 \hat{p} \approx -0.013$,
about $75$ times weaker, because the sigmoid is flat at $z = -5$ and its
derivative $\sigma'(z)$ multiplies everything. Squared error learns
slowest from exactly the examples it gets most wrong. In cross-entropy,
the logarithm cancels the flatness of the sigmoid, and the gradient is
largest when the model is most wrong (\Cref{fig:ch07-ce-vs-se}).

\begin{figure}[ht]
  \centering
  \begin{tikzpicture}
    \begin{axis}[
      name=loss, width=0.48\textwidth, height=0.36\textwidth,
      xmin=0, xmax=1, ymin=0, ymax=4.6,
      xtick={0, 0.5, 1}, ytick={0, 1, 2, 3, 4},
      xlabel={$q_c$, probability of the correct class}, ylabel={loss},
      axis lines*=left, title={\small loss},
      legend style={at={(0.98,0.95)}, anchor=north east, draw=none,
        font=\footnotesize},
      legend cell align=left,
    ]
      \addplot[very thick, blue!60!black, domain=0.01:1, samples=200]
        {-ln(x)};
      \addplot[very thick, dashed, accent, domain=0:1, samples=50]
        {(1-x)^2};
      \legend{cross-entropy, squared error}
    \end{axis}
    \begin{axis}[
      at={(loss.east)}, anchor=west, xshift=0.07\textwidth,
      width=0.48\textwidth, height=0.36\textwidth,
      xmin=-6, xmax=6, ymin=0, ymax=1.05,
      xtick={-6, -3, 0, 3, 6}, ytick={0, 0.5, 1},
      xlabel={logit $z$ (label $y = 1$)}, ylabel={size of gradient},
      axis lines*=left, title={\small gradient},
    ]
      \addplot[very thick, blue!60!black, domain=-6:6, samples=150]
        {1/(1+exp(x))};
      \addplot[very thick, dashed, accent, domain=-6:6, samples=150]
        {2/(1+exp(-x))/(1+exp(x))^2};
    \end{axis}
  \end{tikzpicture}
  \caption{Cross-entropy (solid) and squared error (dashed) for a
    prediction that should be $1$. Left: as a function of the predicted
    probability, cross-entropy grows without bound as the prediction
    approaches $0$, while squared error never exceeds $1$. Right: the size
    of the gradient with respect to the logit of a sigmoid output.
    Cross-entropy pushes hardest on the most wrong predictions; the
    gradient of squared error fades away on them.}
  \label{fig:ch07-ce-vs-se}
\end{figure}

\subsection{Soft Targets}
\label{sec:ch07-soft-targets}

Nothing in the definition requires the target to be one-hot. Any
distribution $\vect{p}$ can be the target, and the loss $H(\vect{p},
\vect{q}) = -\sum_k p_k \log q_k$ has the gradient $\vect{q} - \vect{p}$
with respect to the logits, by the same calculation as
\cref{eq:ch07-ce-grad}. By Gibbs' inequality, it is smallest when the
prediction equals the target, $\vect{q} = \vect{p}$, where it equals the
entropy $H(\vect{p})$, not $0$. Two widely used techniques train on soft
targets.

\paragraph{Label smoothing.}
With one-hot targets, the loss keeps falling as $q_c \to 1$, which
requires the correct logit to run off to infinity. The model is rewarded
for ever more extreme confidence, and it tends to become overconfident.
\emph{Label smoothing}~\cite{szegedy2016rethinking} replaces the target
by a mixture of the one-hot label and the uniform distribution,
\[
  \vect{p} = (1 - \varepsilon)\, \vect{y} + \frac{\varepsilon}{K}\,\vect{1},
\]
for a small $\varepsilon$ such as $0.1$. With $K = 4$ and the first class
correct, the target is $(0.925, 0.025, 0.025, 0.025)$. The best
prediction is now this target itself, with finite logits, and the
minimum loss is its entropy, $0.35$ nats. For the prediction $\vect{q}
\approx (0.61, 0.22, 0.14, 0.03)$ above, the smoothed loss is $0.63$
nats, against $0.50$ with the hard label: the smoothed loss also charges
for putting only $0.03$ on the last class.

\paragraph{Knowledge distillation.}
A large, accurate \emph{teacher} model can train a small \emph{student}
model by using its predicted distributions as the targets
\cite{hinton2015distilling}. A teacher that says ``cat $0.7$, dog
$0.25$, car $0.05$'' conveys more than the label ``cat'': it says that
this cat looks somewhat like a dog and not at all like a car. The student
minimizes the cross-entropy $H(\vect{p}_{\text{teacher}},
\vect{q}_{\text{student}})$, usually with both distributions softened by
a temperature $T > 1$ (\Cref{sec:ch07-entropy-ai}) to bring out the
small probabilities, and usually combined with the ordinary loss on the
true labels.

\subsection{Cross-Entropy and Language Models}
\label{sec:ch07-ce-lm}

A language model is a classifier whose classes are the tokens of its
vocabulary. At each position of a text it predicts a distribution
$\vect{q}_t$ for the next token from the tokens before it, and its loss
is the cross-entropy against the token that actually comes next,
averaged over the $n$ positions:
\[
  \mathcal{L} = -\frac{1}{n} \sum_{t=1}^{n} \log q_t(x_t) .
\]
This is the maximum likelihood objective of \Cref{sec:ch06-mle-ml}, and
the cross-entropy reading adds three things.

\paragraph{A baseline.}
A model that knows nothing and predicts the uniform distribution over a
vocabulary of $V$ tokens has loss $\log V$ on every token. For the
$50{,}257$-token vocabulary of GPT-2, this is $\log 50{,}257 \approx
10.8$ nats. A freshly initialized model predicts nearly uniformly, so
training should start near this value, and a loss that starts far above
it signals a bug, such as a badly scaled
initialization.

\paragraph{Perplexity.}
The perplexity of a model on a text is $e^{\mathcal{L}}$, the exponential
of the cross-entropy. The uniform model has perplexity exactly $V$, and
in general a perplexity of $k$ means the model is as uncertain, on
average, as a fair choice among $k$ tokens. A loss of $2.3$ nats per
token is a perplexity of about $10$; reducing the loss to $2.0$ lowers it
to about $7.4$. Because the scale is exponential, small differences in
loss are large differences in perplexity. Perplexities are only
comparable between models that use the same tokenizer, because a model
with longer tokens makes fewer, harder predictions; to compare across
tokenizers, papers report \emph{bits per byte}, the total loss in bits
divided by the length of the text in bytes.

\paragraph{A floor.}
By Gibbs' inequality, the loss of any model is at least the entropy of
the text itself, the uncertainty that remains in the next token even for
a perfect model. ``The capital of France is'' has one likely next token,
but ``My favorite color is'' has many, and no amount of training makes
the second one predictable. The loss of a language model therefore does
not go to $0$; it approaches this floor, and only the part above it,
$H(p, q) - H(p)$, can be removed by better models. That part is the KL
divergence of the next section.

\begin{aibox}
  Cross-entropy is the most widely used loss in deep learning.
  \begin{itemize}
    \item Image classifiers, speech recognizers, and the pretraining and
      fine-tuning of language models all minimize the cross-entropy
      between the labels or next tokens and the model's softmax output.
    \item By Shannon's source coding theorem, a model with cross-entropy
      $\mathcal{L}$ bits per token can be turned into a compressor that
      stores text in about $\mathcal{L}$ bits per token, using a
      technique called arithmetic coding. Lowering the loss of a language
      model and compressing text better are the same goal, which is why
      language modeling is sometimes described as compression.
    \item Label smoothing with $\varepsilon = 0.1$ is a standard setting
      for training image classifiers and was used to train the original
      transformer~\cite{vaswani2017attention}.
  \end{itemize}
\end{aibox}

\subsection{In Code: Cross-Entropy Losses}
\label{sec:ch07-ce-code}

The definition is one line of NumPy. For the weather station, dividing by
$\log 2$ converts nats to bits:

\begin{lstlisting}
import numpy as np
from scipy import special

def cross_entropy(p, q):
    return -np.sum(p * np.log(q))

summer = np.array([0.5, 0.25, 0.125, 0.125])
winter = np.array([0.125, 0.25, 0.125, 0.5])
cross_entropy(winter, summer) / np.log(2)   # 2.5 bits: wrong code
cross_entropy(winter, winter) / np.log(2)   # 1.75 bits: the entropy
\end{lstlisting}

A classifier's loss should be computed from the logits rather than from
the probabilities. Using $-\log q_c = \log\sum_j e^{z_j} - z_c$ avoids
taking the log of a probability that has underflowed to $0$, and SciPy's
\texttt{logsumexp} subtracts the largest logit first, as the softmax of
\Cref{sec:ch01-ai-functions} does:

\begin{lstlisting}
z = np.array([2.0, 1.0, 0.5, -1.0])
-np.log(special.softmax(z)[0])          # 0.495
special.logsumexp(z) - z[0]             # 0.495: the same

z = np.array([0.0, -800.0])             # correct class: the second
-np.log(special.softmax(z)[1])          # inf: q underflowed to 0
special.logsumexp(z) - z[1]             # 800.0: the true loss
\end{lstlisting}

For a batch of $N$ examples, the logits form an $N \times K$ matrix and
the labels a vector of class indices. Fancy indexing picks out each row's
correct logit:

\begin{lstlisting}
rng = np.random.default_rng(0)
Z = rng.normal(size=(5, 4))             # logits: 5 examples, 4 classes
y = np.array([0, 3, 1, 1, 2])           # correct classes
losses = special.logsumexp(Z, axis=1) - Z[np.arange(5), y]
losses.mean()                           # 1.366
\end{lstlisting}

PyTorch~\cite{paszke2019pytorch} provides the same computation, with label
smoothing and soft targets built in, and its automatic differentiation
confirms the gradient $\vect{q} - \vect{y}$ of \cref{eq:ch07-ce-grad}:

\begin{lstlisting}
import torch
import torch.nn.functional as F

Zt, yt = torch.tensor(Z, dtype=torch.float32), torch.tensor(y)
F.cross_entropy(Zt, yt)                       # 1.366: same as above
F.cross_entropy(Zt, yt, label_smoothing=0.1)  # 1.384

z = torch.tensor([2.0, 1.0, 0.5, -1.0], requires_grad=True)
loss = F.cross_entropy(z.unsqueeze(0), torch.tensor([0]))
loss.backward()
z.grad                  # [-0.391  0.224  0.136  0.030] = q - y
\end{lstlisting}

The target can also be a matrix of probabilities, one row per example,
for distillation or hand-made smoothing. For two classes,
\texttt{F.binary\_cross\_entropy\_with\_logits} takes the single logit
$z$ of each example and computes the binary cross-entropy of
$\sigma(z)$ stably.

\begin{warning}
  \texttt{F.cross\_entropy} and \texttt{nn.CrossEntropyLoss} expect
  \emph{logits} and apply the softmax themselves. Passing probabilities
  applies it twice, which squashes the predictions toward uniform. The
  code still runs and the loss still falls, but it can never reach
  $0$: for the logits $(10, 0, 0, 0)$ and the correct first class, the
  loss is $0.0001$, but on the softmax output it is $0.74$, and the
  gradients are distorted accordingly.
\end{warning}

\section{KL Divergence}
\label{sec:ch07-kl}

\subsection{The Extra Cost of the Wrong Model}
\label{sec:ch07-kl-def}

Gibbs' inequality says that describing data from $p$ with a model $q$
costs $H(p, q)$, never less than the $H(p)$ that a model of $p$ itself
would need. The entropy $H(p)$ belongs to the data and no model can
remove it. The excess belongs to the model: it is how much worse $q$ is
than the truth. That excess has a name.

\begin{definition}[KL divergence]
  \label{def:kl}
  The \emph{Kullback--Leibler divergence}, or \emph{KL divergence}, of a
  distribution $q$ from a distribution $p$ is
  \[
    D_{\mathrm{KL}}(p \,\|\, q) = H(p, q) - H(p)
    = \sum_x p(x) \log \frac{p(x)}{q(x)}
    = \E_{X \sim p}\Bigl[\log \frac{p(X)}{q(X)}\Bigr] ,
  \]
  with terms where $p(x) = 0$ taken to be $0$. It is infinite if $q(x)
  = 0$ for some $x$ with $p(x) > 0$. For densities, the sum is an
  integral.
\end{definition}

The middle formula follows by combining the two sums: $-\sum_x p(x) \log
q(x) + \sum_x p(x) \log p(x) = \sum_x p(x) \log\bigl(p(x)/q(x)\bigr)$.
Each term compares the probability the truth gives an outcome with the
probability the model gives it, through the log-ratio $\log(p/q)$, and
weights the comparison by how often the outcome occurs. For the weather
station of \Cref{sec:ch07-ce-idea},
\[
  D_{\mathrm{KL}}(p_{\text{winter}} \,\|\, p_{\text{summer}})
  = 2.5 - 1.75 = 0.75 \text{ bits},
\]
the cost of sending winter reports with the summer code. Gibbs'
inequality, restated, is the basic property of the KL divergence.

\begin{proposition}[KL divergence is nonnegative]
  \label{prop:ch07-kl-nonneg}
  $D_{\mathrm{KL}}(p \,\|\, q) \ge 0$, with equality only when $q = p$.
\end{proposition}

So the KL divergence behaves like a measure of how far $q$ is from $p$:
it is $0$ when they agree and positive when they do not. It is not a
distance in the sense of geometry, though, in two ways.

\paragraph{It is not symmetric.}
The truth and the model play different roles, and swapping them changes
the answer. Take a fair coin $p = (0.5, 0.5)$ and a biased coin $q = (0.9,
0.1)$:
\[
  D_{\mathrm{KL}}(p \,\|\, q)
  = 0.5 \log_2 \frac{0.5}{0.9} + 0.5 \log_2 \frac{0.5}{0.1}
  \approx 0.74 \text{ bits},
  \qquad
  D_{\mathrm{KL}}(q \,\|\, p) \approx 0.53 \text{ bits}.
\]
Modeling a fair coin as biased is worse than the reverse, because the
biased model is badly surprised by the many tails a fair coin produces.
The double bar $\|$ is a reminder that the two arguments are not
interchangeable; this is a \emph{divergence}, not a distance. In the
extreme case, a model that rules out something the truth allows has
infinite divergence, while the reverse is finite: $D_{\mathrm{KL}}(p \,\|\,
q) = \infty$ for the fair coin and a two-headed model, but
$D_{\mathrm{KL}}(q \,\|\, p) = \log 2$.

\paragraph{It does not obey the triangle inequality.}
For a distance, going from $p$ to $r$ directly is never longer than going
through a third point $q$. With coins of probability of heads $0.5$,
$0.7$, and $0.9$, the KL divergence from the first to the third is
$0.74$ bits, but from the first to the second plus from the second to the
third is only $0.35$ bits.

Neither defect matters much in practice. What matters is that the KL
divergence is $0$ exactly when the model is right, positive otherwise,
and, as the next subsection shows, that it is what training minimizes.

Unlike differential entropy (\Cref{sec:ch07-differential}), the KL
divergence between densities does not depend on units. Changing the
units of $x$ stretches both densities by the same factor, which cancels
in the ratio $p/q$. It is also never negative, so for continuous
distributions it is the quantity with a clean meaning, and it is the one
used in practice.

\subsection{Minimizing KL Divergence Is Maximum Likelihood}
\label{sec:ch07-kl-mle}

Suppose the data come from an unknown distribution $p_{\text{data}}$ and
we fit a model $q_\theta$. The ideal goal is to make the model as close
to the truth as possible, that is, to minimize $D_{\mathrm{KL}}(p_{\text{data}}
\,\|\, q_\theta)$. Splitting the divergence,
\[
  D_{\mathrm{KL}}(p_{\text{data}} \,\|\, q_\theta)
  = \underbrace{H(p_{\text{data}}, q_\theta)}_{\text{depends on } \theta}
  - \underbrace{H(p_{\text{data}})}_{\text{does not}} .
\]
The entropy of the data is a fixed number, unknown but untouched by
$\theta$, so minimizing the KL divergence and minimizing the
cross-entropy have the same solution. We cannot compute either exactly,
since $p_{\text{data}}$ is unknown, but we have samples $x_1, \dots, x_N$
from it, and the cross-entropy is an expectation under $p_{\text{data}}$,
which a sample average estimates:
\[
  H(p_{\text{data}}, q_\theta)
  = \E_{X \sim p_{\text{data}}}\bigl[-\log q_\theta(X)\bigr]
  \approx -\frac{1}{N} \sum_{i=1}^{N} \log q_\theta(x_i)
  = -\frac{1}{N}\, \ell(\theta) .
\]
This is the average negative log-likelihood of \Cref{chap:ch06}. Three
descriptions of training therefore coincide:
\[
  \text{maximize the likelihood}
  \;\Longleftrightarrow\;
  \text{minimize the cross-entropy}
  \;\Longleftrightarrow\;
  \text{minimize } D_{\mathrm{KL}}(p_{\text{data}} \,\|\, q_\theta) .
\]
The last one says what maximum likelihood is aiming for: the member of
the model family closest to the truth, in the sense of the KL divergence.
It also explains the floor of \Cref{sec:ch07-ce-lm}. The training loss of
a language model is $H(p_{\text{data}}) + D_{\mathrm{KL}}(p_{\text{data}}
\,\|\, q_\theta)$; better models shrink the divergence, and the entropy
of language is what remains.

\subsection{Forward and Reverse KL}
\label{sec:ch07-kl-direction}

Since the KL divergence is not symmetric, fitting a model $q$ to a
distribution $p$ can mean two different things, and the choice changes
the answer. Minimizing the \emph{forward} KL, $D_{\mathrm{KL}}(p \,\|\,
q)$, is what maximum likelihood does. Minimizing the \emph{reverse} KL,
$D_{\mathrm{KL}}(q \,\|\, p)$, is what several methods do when $p$ can be
evaluated but not sampled, as we will see below. The difference is
clearest when the model is too simple to match the truth.

Let $p$ be an equal mixture of two Gaussians, centered at $-2$ and $2$
with standard deviation $0.6$, and fit a single Gaussian $q$ to it.
\begin{itemize}
  \item \emph{Forward KL} averages $\log(p/q)$ over $p$. Wherever $p$ is
    large and $q$ is small, the ratio is huge, so the forward KL forces
    $q$ to put probability everywhere $p$ does. The best Gaussian matches
    the mean and variance of $p$: it is $\mathcal{N}(0, 2.09^2)$, wide
    enough to cover both bumps, at the price of putting much of its mass
    in the middle, where $p$ has almost none. This behavior is called
    \emph{mode covering}.
  \item \emph{Reverse KL} averages $\log(q/p)$ over $q$. Now the penalty
    falls where $q$ is large and $p$ is small, so $q$ must avoid every
    region that $p$ considers unlikely, and nothing forces it to cover all
    of $p$. The best Gaussian settles on one bump, $\mathcal{N}(2.00,
    0.60^2)$ (or its mirror image at $-2$), and ignores the other. This is
    \emph{mode seeking}. Its divergence is about $0.69 \approx \log 2$
    nats, the cost of missing half of $p$.
\end{itemize}
\Cref{fig:ch07-kl-direction} shows both fits. Neither is wrong; they
answer different questions. The forward fit never misses anything that
can happen but produces unlikely samples from the gap. The reverse fit
produces only realistic samples but never produces the left bump.

\begin{figure}[ht]
  \centering
  \begin{tikzpicture}
    \begin{axis}[
      width=0.75\textwidth, height=0.36\textwidth,
      xmin=-6, xmax=6, ymin=0, ymax=0.72,
      xtick={-4, -2, 0, 2, 4}, ytick={0, 0.2, 0.4, 0.6},
      xlabel={$x$}, ylabel={density},
      axis lines*=left,
      legend style={at={(0.02,0.97)}, anchor=north west, draw=none,
        font=\footnotesize},
      legend cell align=left,
    ]
      \addplot[fill=gray!25, draw=gray!60, domain=-6:6, samples=200]
        {0.5*exp(-(x+2)^2/0.72)/(0.6*sqrt(2*pi))
         + 0.5*exp(-(x-2)^2/0.72)/(0.6*sqrt(2*pi))} \closedcycle;
      \addplot[very thick, blue!60!black, domain=-6:6, samples=150]
        {exp(-x^2/(2*4.36))/(2.088*sqrt(2*pi))};
      \addplot[very thick, dashed, accent, domain=-6:6, samples=150]
        {exp(-(x-2)^2/(2*0.366))/(0.605*sqrt(2*pi))};
      \legend{$p$ (two bumps), {forward KL fit, $\mathcal{N}(0, 2.09^2)$},
        {reverse KL fit, $\mathcal{N}(2, 0.60^2)$}}
    \end{axis}
  \end{tikzpicture}
  \caption{Fitting one Gaussian to a distribution with two bumps (gray).
    Minimizing $D_{\mathrm{KL}}(p \,\|\, q)$ gives a wide Gaussian that
    covers both bumps (solid). Minimizing $D_{\mathrm{KL}}(q \,\|\, p)$
    gives a narrow Gaussian that sits on one bump (dashed).}
  \label{fig:ch07-kl-direction}
\end{figure}

\subsection{KL Divergence Between Gaussians}
\label{sec:ch07-kl-gaussian}

For two Gaussians, the KL divergence has a closed form, which is why it
appears so often in models built from Gaussians. Let $p =
\mathcal{N}(\mu_1, \sigma_1^2)$ and $q = \mathcal{N}(\mu_2, \sigma_2^2)$.
From the Gaussian log-density,
\[
  \log \frac{p(x)}{q(x)}
  = \log \frac{\sigma_2}{\sigma_1}
  - \frac{(x - \mu_1)^2}{2\sigma_1^2} + \frac{(x - \mu_2)^2}{2\sigma_2^2} .
\]
Under $p$, the mean of $(x - \mu_1)^2$ is $\sigma_1^2$, and the mean of
$(x - \mu_2)^2$ is $\sigma_1^2 + (\mu_1 - \mu_2)^2$, the variance plus the
squared offset of the mean. So
\begin{equation}
  D_{\mathrm{KL}}(p \,\|\, q)
  = \log \frac{\sigma_2}{\sigma_1}
  + \frac{\sigma_1^2 + (\mu_1 - \mu_2)^2}{2\sigma_2^2} - \frac12 .
  \label{eq:ch07-kl-gauss}
\end{equation}
With equal variances it reduces to $(\mu_1 - \mu_2)^2/(2\sigma^2)$, half
the squared distance between the means in units of the standard
deviation: $D_{\mathrm{KL}}(\mathcal{N}(1, 1) \,\|\, \mathcal{N}(0, 1)) =
\tfrac12$. With equal means, it measures the mismatch in spread, and it
is not symmetric: $D_{\mathrm{KL}}(\mathcal{N}(0, 1) \,\|\, \mathcal{N}(0,
2^2)) \approx 0.32$, while $D_{\mathrm{KL}}(\mathcal{N}(0, 2^2) \,\|\,
\mathcal{N}(0, 1)) \approx 0.81$. A model that is too narrow is penalized
more than one that is too wide, as in the forward fit above.

The case used most is a Gaussian measured against the standard Gaussian.
With $\mu_2 = 0$ and $\sigma_2 = 1$, \cref{eq:ch07-kl-gauss} becomes
\[
  D_{\mathrm{KL}}\bigl(\mathcal{N}(\mu, \sigma^2) \,\|\, \mathcal{N}(0, 1)\bigr)
  = \tfrac12 \bigl( \mu^2 + \sigma^2 - 1 - \log \sigma^2 \bigr) .
\]
It is $0$ only at $\mu = 0$, $\sigma = 1$, and it grows as the mean moves
away from $0$ or the variance away from $1$. For $\mu = 0.5$ and $\sigma
= 0.8$, it is $0.17$ nats. For multivariate Gaussians with diagonal
covariance, the coordinates are independent, and the divergence is the
sum of this expression over the coordinates. For general covariances,
\[
  D_{\mathrm{KL}}\bigl(\mathcal{N}(\vect{\mu}_1, \Sigma_1) \,\|\,
    \mathcal{N}(\vect{\mu}_2, \Sigma_2)\bigr)
  = \tfrac12 \Bigl[ \operatorname{tr}(\Sigma_2^{-1}\Sigma_1)
    + (\vect{\mu}_2 - \vect{\mu}_1)^\top \Sigma_2^{-1}
      (\vect{\mu}_2 - \vect{\mu}_1)
    - d + \log \frac{\det \Sigma_2}{\det \Sigma_1} \Bigr] ,
\]
which reduces to \cref{eq:ch07-kl-gauss} for $d = 1$.

\subsection{KL Divergence in AI}
\label{sec:ch07-kl-ai}

Besides hiding inside every cross-entropy loss, the KL divergence
appears explicitly wherever a model must stay close to some reference
distribution.

\paragraph{Variational autoencoders.}
A variational autoencoder~\cite{kingma2014vae} encodes each input
$\vect{x}$ as a Gaussian over a latent space, with a mean
$\vect{\mu}(\vect{x})$ and a variance $\sigma_j^2(\vect{x})$ for each
coordinate, samples a code from it, and decodes the code back into an
input. Its loss
has two terms: a reconstruction error, and the KL divergence from each
encoding to the standard Gaussian, computed coordinate by coordinate with
the closed form of the last subsection. The KL term keeps the encodings
close to $\mathcal{N}(\vect{0}, I)$, so that codes drawn from the
standard Gaussian later decode into realistic new inputs. Mathematically,
the VAE fits an approximate posterior by minimizing a reverse KL, since
the true posterior can be evaluated, up to a constant, but not sampled.

\paragraph{Fine-tuning with a leash.}
When a language model is fine-tuned with reinforcement learning from
human feedback~\cite{ouyang2022instructgpt}, the policy $\pi$ is
trained to earn high reward $r$ from a learned reward model, but a
penalty keeps it close to the reference model $\pi_{\text{ref}}$ it
started from:
\[
  \text{maximize}\quad
  \E\bigl[r\bigr] - \beta\, D_{\mathrm{KL}}(\pi \,\|\, \pi_{\text{ref}}) .
\]
Without the penalty, the policy finds outputs that fool the reward model
and drifts into text that no longer reads naturally. The strength $\beta$
trades reward against staying close, in the same way that $\lambda$
trades fit against a prior in \Cref{sec:ch06-map-ml}.

\paragraph{Distillation and soft targets.}
The distillation loss of \Cref{sec:ch07-soft-targets} is often written
as $D_{\mathrm{KL}}(\vect{p}_{\text{teacher}} \,\|\,
\vect{q}_{\text{student}})$. It differs from the cross-entropy only by the
teacher's entropy, which does not depend on the student, so the two give
the same gradients. The KL form has the advantage of being $0$ when the
student matches the teacher exactly.

\paragraph{Mutual information.}
\Cref{sec:ch07-entropy-props} described the mutual information between
$X$ and $Y$ as the shortfall $H(X) + H(Y) - H(X, Y)$. It is also a KL
divergence,
\[
  I(X; Y) = D_{\mathrm{KL}}\bigl(p(x, y) \,\|\, p(x)\,p(y)\bigr),
\]
the distance of the joint distribution from the one that would hold if
$X$ and $Y$ were independent. By \Cref{prop:ch07-kl-nonneg} it is never
negative, and it is $0$ exactly when $X$ and $Y$ are independent.
Contrastive learning methods, which train image and text encoders to
recognize matching pairs, can be read as maximizing a bound on the mutual
information between the two views.

\begin{aibox}
  When you meet a KL term in a paper, ask two questions. Which
  distribution is first? Forward KL, with the data first, covers
  everything the data does; reverse KL, with the model first, picks a
  mode and stays inside it. And what is being held close to what? In VAEs
  the encodings are held near a standard Gaussian; in RLHF the policy is
  held near the pretrained model; in t-SNE~\cite{vandermaaten2008tsne}
  the neighbor probabilities of the low-dimensional map are held near
  those of the original data.
\end{aibox}

\subsection{In Code: Computing KL Divergences}
\label{sec:ch07-kl-code}

SciPy's \texttt{stats.entropy} computes the KL divergence when given two
distributions, and \texttt{special.rel\_entr} gives the individual terms
$p \log(p/q)$:

\begin{lstlisting}
import numpy as np
from scipy import special, stats

summer = np.array([0.5, 0.25, 0.125, 0.125])
winter = np.array([0.125, 0.25, 0.125, 0.5])
stats.entropy(winter, summer, base=2)             # 0.75 bits
special.rel_entr(winter, summer).sum() / np.log(2)  # 0.75: same

cross = -np.sum(winter * np.log2(summer))         # 2.5
cross - stats.entropy(winter, base=2)             # 0.75 = H(p,q) - H(p)

stats.entropy([0.5, 0.5], [0.9, 0.1], base=2)     # 0.737
stats.entropy([0.9, 0.1], [0.5, 0.5], base=2)     # 0.531: not symmetric
stats.entropy([0.5, 0.5], [1.0, 0.0])             # inf
\end{lstlisting}

When no closed form exists, the KL divergence is an expectation under
$p$, so it can be estimated by averaging $\log p(x) - \log q(x)$ over
samples from $p$. Testing this on two Gaussians, where
\cref{eq:ch07-kl-gauss} gives the exact answer $\tfrac12$:

\begin{lstlisting}
rng = np.random.default_rng(0)
x = rng.normal(loc=1.0, scale=1.0, size=100_000)    # samples from p
np.mean(stats.norm.logpdf(x, 1.0, 1.0)
        - stats.norm.logpdf(x, 0.0, 1.0))           # 0.499
\end{lstlisting}

This Monte Carlo estimate is how the KL penalty of RLHF is computed in
practice: the log-probabilities of sampled text under the policy and
under the reference model are subtracted and averaged. PyTorch has the
closed forms for common distributions:

\begin{lstlisting}
import torch
import torch.nn.functional as F
from torch.distributions import Normal, kl_divergence

kl_divergence(Normal(0.5, 0.8), Normal(0.0, 1.0))   # 0.168: the VAE term
kl_divergence(Normal(1.0, 1.0), Normal(0.0, 1.0))   # 0.5

p = torch.tensor([[0.5, 0.5]])                  # one row per example
q = torch.tensor([[0.9, 0.1]])
F.kl_div(q.log(), p, reduction="batchmean")     # 0.511 = KL(p || q)
\end{lstlisting}

\begin{warning}
  \texttt{F.kl\_div(input, target)} computes $D_{\mathrm{KL}}(\text{target}
  \,\|\, \text{input})$, the reverse of the order in which the arguments
  are written, and it expects \texttt{input} as \emph{log}-probabilities.
  Swapping the arguments silently computes the other divergence ($0.368$
  instead of $0.511$ above), and passing probabilities instead of
  log-probabilities gives meaningless numbers, here $-1.19$. Use
  \texttt{reduction="batchmean"}: the default, \texttt{"mean"}, also
  divides by the number of classes and halves the answer for two classes.
\end{warning}

\section{Summary}
\label{sec:ch07-summary}

\begin{itemize}
  \item The \emph{surprise} of an outcome of probability $p$ is $-\log p$:
    rarer outcomes carry more information, a certain outcome carries
    none, and the surprises of independent events add. With $\log_2$ it
    is measured in bits, with the natural logarithm in nats.
  \item The \emph{entropy} $H(X) = -\sum_x p(x) \log p(x)$ is the average
    surprise, a measure of how uncertain $X$ is. A fair coin has $1$ bit
    and a coin with $\theta = 0.9$ has $0.47$. By Shannon's source coding
    theorem, it is also the fewest bits per outcome that any code can
    use on average.
  \item $0 \le H(X) \le \log K$ for $K$ outcomes, with $0$ only when one
    outcome is certain and $\log K$ only for the uniform distribution.
    Entropy is concave, and the entropies of independent variables add.
  \item The \emph{differential entropy} of a density is not a count of
    bits: it can be negative, and it depends on units. The Gaussian
    $\mathcal{N}(\mu, \sigma^2)$ has $h = \tfrac12 \log(2\pi e \sigma^2)$,
    the largest of any density with that variance.
  \item In AI, the entropy of a softmax output measures a model's
    uncertainty, and the temperature controls it. Entropy also drives
    decision-tree splits and exploration bonuses in reinforcement
    learning. In code, use \texttt{scipy.stats.entropy} or
    \texttt{special.entr}, which handle $0 \log 0 = 0$.
  \item The \emph{cross-entropy} $H(p, q) = -\sum_x p(x) \log q(x)$ is the
    average surprise of a model $q$ at data from $p$: the cost of a code
    built for $q$. By \emph{Gibbs' inequality}, $H(p, q) \ge H(p)$, with
    equality only for $q = p$. It is not symmetric, and it is infinite
    if $q$ rules out an outcome that $p$ allows.
  \item With a one-hot label, the cross-entropy is $-\log q_c$, the
    negative log-likelihood of the correct class, so cross-entropy
    training is maximum likelihood. Its gradient with respect to the
    logits is $\vect{q} - \vect{y}$, which stays large for confidently
    wrong predictions, where squared error on a sigmoid output stalls.
  \item Soft targets generalize the loss: label smoothing mixes the label
    with the uniform distribution to curb overconfidence, and
    distillation uses a teacher's predictions as targets. A language
    model's loss starts near $\log V$, its perplexity is
    $e^{\text{loss}}$, and its loss cannot go below the entropy of the
    text. In code, compute the loss from logits with log-sum-exp, and
    pass logits, not probabilities, to \texttt{F.cross\_entropy}.
  \item The \emph{KL divergence} $D_{\mathrm{KL}}(p \,\|\, q) = H(p, q) -
    H(p) = \sum_x p(x) \log\bigl(p(x)/q(x)\bigr)$ is the extra cost of
    the wrong model. It is never negative and is $0$ only when $q = p$,
    but it is not a distance: it is not symmetric and breaks the triangle
    inequality. Unlike differential entropy, it does not depend on units.
  \item Maximizing likelihood, minimizing cross-entropy, and minimizing
    $D_{\mathrm{KL}}(p_{\text{data}} \,\|\, q_\theta)$ are the same goal,
    because $H(p_{\text{data}})$ does not depend on the model.
  \item Forward KL, $D_{\mathrm{KL}}(p \,\|\, q)$, makes the model cover
    everything the data does; reverse KL, $D_{\mathrm{KL}}(q \,\|\, p)$,
    makes it settle on one mode. Between Gaussians the KL divergence has
    a closed form, and against $\mathcal{N}(0, 1)$ it is $\tfrac12(\mu^2
    + \sigma^2 - 1 - \log \sigma^2)$, the regularizer of a VAE.
  \item KL terms keep a model near a reference: VAE encodings near a
    standard Gaussian, an RLHF policy near its pretrained model, a student
    near its teacher. The mutual information $I(X; Y)$ is the KL
    divergence from independence. In code, \texttt{F.kl\_div(input,
    target)} computes $D_{\mathrm{KL}}(\text{target} \,\|\, \text{input})$
    and expects log-probabilities as \texttt{input}.
\end{itemize}

Cover and Thomas~\cite{cover2006elements} is the standard introduction to
information theory, including the source coding theorem and mutual
information, which this chapter only stated. Shannon's original
paper~\cite{shannon1948mathematical} is still readable today.
